\begin{quote}
\textbf{Maturity: {[}Verified{]}.} The February 2026 draft of this paper
declined to publish figures it could not measure. Most of what it
labelled \emph{Specified} or \emph{Planned} is now live on the chain-id
40204 testnet. This revision re-anchors every mechanical claim to a file
path or a deployed contract address, and it keeps the earlier draft's
discipline: where a number is a documented target rather than a measured
result, it is labelled as one.
\end{quote}

\subsection{Abstract}\label{abstract}

We present the protocol specification for the Citrate Network, a
BlockDAG system that extends GhostDAG consensus with AI-native
execution. The architecture comprises three layers: a \textbf{GhostDAG
consensus layer} providing parallel block production with committee BFT
finality checkpoints (k=18, up to 10 parents per block, a 100-validator
committee with a 67-signature quorum); a \textbf{Lattice Virtual Machine
(LVM)} providing EVM bytecode compatibility via REVM, augmented with
families of AI-specific precompiled contracts for inference, verifiable
proof checking, deterministic fixed-point tensor compute, and
paraconsistent learning aggregation; and a \textbf{Model Context
Protocol (MCP)} layer providing OpenAI- and Anthropic-compatible REST
endpoints for model discovery, inference, and orchestration. This paper
documents the implemented components, the consensus engine, the virtual
machine with tested AI precompiles, dual ECDSA/Ed25519 identity, and
SALT tokenomics realized in deployed contracts, and distinguishes them
from components that remain feature-gated or in design. The consensus
layer is integrated with the federated learning protocol described in
companion Paper II, whose Belnap aggregation precompile is live. The
native token SALT (1 billion supply) functions as gas, staking
collateral, and governance weight. As of this revision the testnet
carries \textbf{39 deployed contracts} on chain 40204 and the node's
core workspace holds \textbf{3,517 tests}.

\textbf{Keywords:} BlockDAG, GhostDAG, BFT finality, EVM compatibility,
AI precompiles, Model Context Protocol, verifiable inference, Citrate
Network

\subsection{1. Introduction}\label{introduction}

Citrate is a three-layer blockchain designed for AI-native applications.
The thesis is that a general-purpose execution environment with on-chain
AI primitives, model registration, verifiable inference, deterministic
tensor operations, adapter management, and model orchestration, can
serve as infrastructure for decentralized AI in a way that purpose-built
AI coordination networks (which lack general programmability) and
general-purpose blockchains (which lack AI-specific optimizations)
cannot.

The design draws on three established technologies. First,
\textbf{GhostDAG consensus}, originated by Sompolinsky and Zohar {[}1,
2{]} and first productionized by Kaspa, provides a BlockDAG structure
enabling parallel block production with deterministic ordering. Second,
the \textbf{Ethereum Virtual Machine}, the most widely deployed
smart-contract runtime, provides developer tooling, auditing
infrastructure, and application compatibility. Third, the \textbf{Model
Context Protocol}, originated by Anthropic {[}23{]}, provides
standardized interfaces for model-to-model and model-to-tool
communication.

Since the February 2026 draft, a category of AI-native and
inference-verifying blockchains has matured, and this specification
positions Citrate against it explicitly (Section 7): the systematization
of decentralized AI (DeAI) by Zhang et al.~{[}31{]}; proof-of-logits and
subnet-scored designs (Ambient, Bittensor's Yuma consensus) {[}13{]};
and optimistic and TEE-attested inference rollups {[}32{]}. Citrate's
distinguishing choice is to embed AI operations as composable
precompiles inside a general EVM, rather than as an isolated task
market.

This paper is organized as a protocol specification. Section 2 describes
the consensus layer; Section 3 the Lattice Virtual Machine and its AI
precompiles; Section 4 the MCP orchestration layer; Section 5 SALT
tokenomics; Section 6 security; Section 7 comparative position; Section
8 open problems.

\textbf{Implementation status.} Claims marked \textbf{{[}Implemented{]}}
describe components built and tested on the chain-40204 testnet, with a
file path or contract address given. Claims marked
\textbf{{[}Specified{]}} are designed in detail but not yet built.
Claims marked \textbf{{[}Feature-gated{]}} exist in code but are
compiled out of the default build. This convention lets a reader
distinguish what exists from what is proposed, and check each claim
against the source.

\subsection{2. Consensus Layer: GhostDAG with BFT
Finality}\label{consensus-layer-ghostdag-with-bft-finality}

\subsubsection{2.1 GhostDAG Overview}\label{ghostdag-overview}

\textbf{{[}Implemented{]}} Citrate's consensus layer implements the
GhostDAG protocol {[}1, 2{]}, which generalizes Nakamoto consensus to
directed acyclic graphs (\texttt{core/consensus/src/ghostdag.rs},
\texttt{GhostDag} at line 54). Blocks reference multiple parents (up to
10 in Citrate: one selected parent plus up to nine merge parents;
\texttt{core/consensus/src/types.rs:179}), forming a DAG rather than a
chain. The protocol partitions blocks into a blue set (consistent with
the k-cluster rule, anticone size ≤ k) and a red set. A selected-parent
chain through the blue set provides deterministic total ordering for
contract execution (\texttt{tip\_selection.rs}, \texttt{ordering.rs},
\texttt{chain\_selection.rs}), and blue scores serve as the primary
tip-selection weight (\texttt{calculate\_blue\_set} at
\texttt{ghostdag.rs:201}).

The \textbf{k parameter is 18}
(\texttt{core/consensus/src/types.rs:178}), calibrated for the network's
block time and expected propagation delay. For context, Kaspa's
Crescendo hardfork (May 2025) raised k to 124 at 10 blocks per second
with max parents 16 {[}30{]}. Citrate's more conservative k=18 reflects
a slower block rate chosen to accommodate the larger blocks required by
AI-augmented blocks carrying embedding data (Section 2.3). The k-cluster
rule is exercised by roughly twenty consensus test files, including
\texttt{k\_cluster\_tests.rs}, \texttt{fork\_choice\_tests.rs},
\texttt{h07\_reachability.rs}, and \texttt{proptest\_consensus.rs}.

\subsubsection{2.2 Block Time and Rate}\label{block-time-and-rate}

\textbf{{[}Implemented{]}} The chain-40204 testnet targets a
\textbf{2-second block time} (\texttt{testnet-config.toml:4};
\texttt{core/sequencer/src/block\_builder.rs:76},
\texttt{block\_time\_target:\ 2}); the devnet runs at 1 second. This is
more conservative than the February 2026 draft, which described a
0.5-second target; the code and the running testnet are the ground
truth, and this revision corrects the figure. The design rationale is
unchanged: AI-augmented blocks are larger than pure-transaction blocks.
When the federated learning path of Paper II is active, each block
carries embedding vectors alongside transactions; a 768-dimensional
float32 embedding adds roughly 3 KB per block, comfortably within the
maximum block size at a 2-second cadence. Block time remains a
governance-adjustable parameter, so the rate can be raised if embedding
compression advances or the learning path is not needed on a given
deployment.

\subsubsection{2.3 BFT Finality
Checkpoints}\label{bft-finality-checkpoints}

\textbf{{[}Implemented{]}} Citrate runs a dual finality mechanism.
Depth-based optimistic confirmation is provided by the
\texttt{FinalityTracker} (\texttt{core/consensus/src/finality.rs:95}).
On top of it, \textbf{committee BFT checkpoints}
(\texttt{core/consensus/src/checkpoint.rs}, ``WP-S.3: Committee BFT
Checkpoints'') select a committee deterministically by VRF seed and
validator public key, and require a \textbf{67-of-100 quorum} (2/3+1;
\texttt{checkpoint.rs:88,103}) over a domain-separated ed25519 vote
(\texttt{CITRATE-CHECKPOINT-V1}, \texttt{checkpoint.rs:59}). The default
checkpoint interval is \textbf{50 blocks} (\texttt{checkpoint.rs:101});
when a checkpoint reaches quorum it overrides depth-based finality up to
the checkpoint height, and the committed block hash, blue set, and state
root become irreversible. The README documents a ≤12-second optimistic
finality target; the checkpoint cadence itself is 50 blocks, and this
revision states both rather than collapsing them into a single headline
number. The checkpoint commitment is extensible to learning state
(routing weights, adapter registries), which is how Paper II's learning
checkpoints bind to finality
(\texttt{core/learning/src/checkpoint.rs:14}).

Safety holds if fewer than n/3 committee members are Byzantine, by
standard BFT arguments {[}11{]}; with a 100-validator committee
requiring 67 signatures, this tolerates up to 33 Byzantine validators.
Liveness holds under GhostDAG's blue-set property when fewer than n/2 of
validators are Byzantine {[}1, 2{]}.

\subsubsection{2.4 Performance: What Is and Is Not
Measured}\label{performance-what-is-and-is-not-measured}

The February 2026 draft declined to publish a throughput figure without
a measurement methodology. That discipline is retained, and can now be
stated precisely.

\textbf{{[}Implemented, measured{]}} The parallel executor is
benchmarked with a gate that asserts an \textbf{8-worker throughput of
at least 2× the 1-worker throughput} on disjoint-sender workloads
(\texttt{core/execution/benches/tps\_parallel.rs:10}). This relative
speedup is the claim the in-repo benchmark actually proves. The
executor's absolute ceiling on disjoint transfers is high (prior sprints
measured hundreds of thousands of transfers per second at the executor
layer), but the executor ceiling is not the network TPS.

\textbf{{[}Documented target, not benched{]}} The README advertises
5,000 TPS sustained with a 10,000 TPS ceiling (\texttt{README.md:65}).
This is a documented target, not a value produced by the in-repo
benchmarks, and the benchmark header itself disclaims its absolute
throughput as executor-ceiling-bound (\texttt{tps\_parallel.rs:145}). We
report it as a target and label it as such. Real-world network
throughput is bounded by signature recovery, RLP decode, mempool gossip,
and disk I/O; those are the optimization targets for subsequent work.

\subsection{3. Execution Layer: The Lattice Virtual
Machine}\label{execution-layer-the-lattice-virtual-machine}

\subsubsection{3.1 EVM Compatibility}\label{evm-compatibility}

\textbf{{[}Implemented{]}} The Lattice Virtual Machine (LVM) executes
EVM bytecode via REVM (\texttt{core/execution/src/revm\_adapter.rs:9}),
registering Citrate's custom precompiles through
\texttt{register\_citrate\_precompiles}. Solidity, Vyper, or Yul
compiled for Ethereum deploys on Citrate without recompilation, giving
immediate access to Hardhat, Foundry, OpenZeppelin, Ethers.js, and
existing audit tooling. Parallel execution with MVCC state is
implemented in \texttt{core/execution/src/parallel/} and
\texttt{core/execution/src/mvcc/}, following the precedent of
EVM-compatible chains (Avalanche C-Chain, Moonbeam, Hedera) that pair
EVM compatibility with non-Ethereum consensus.

Each account maintains a \textbf{dual cryptographic identity}: ECDSA
over secp256k1 for Ethereum wallet compatibility and Ed25519 for native
operations, following Hedera's production-proven multi-scheme pattern.
Address handling accepts both embedded 20-byte EVM addresses and full
32-byte public keys, hashing the latter with Keccak-256
(\texttt{core/execution/src/types.rs}).

\subsubsection{3.2 AI Precompiled
Contracts}\label{ai-precompiled-contracts}

\textbf{{[}Implemented{]}} The LVM extends the EVM with several families
of AI-specific precompiles
(\texttt{core/execution/src/precompiles/mod.rs:82-128}). Precompiles
execute native code rather than EVM bytecode, enabling operations
prohibitively expensive as Solidity, following Ethereum's own precompile
pattern (ecrecover at 0x01, and so on, which Citrate also provides at
0x01--0x09). The February 2026 draft anticipated five precompiles at
0x1000--0x1004; the implemented address map is richer and is documented
here as it actually exists.

\textbf{Table 1. Implemented AI precompile families.}

{\def\LTcaptype{none} % do not increment counter
\begin{longtable}[]{@{}
  >{\raggedright\arraybackslash}p{(\linewidth - 6\tabcolsep) * \real{0.2586}}
  >{\raggedright\arraybackslash}p{(\linewidth - 6\tabcolsep) * \real{0.1379}}
  >{\raggedright\arraybackslash}p{(\linewidth - 6\tabcolsep) * \real{0.4655}}
  >{\raggedright\arraybackslash}p{(\linewidth - 6\tabcolsep) * \real{0.1379}}@{}}
\toprule\noalign{}
\begin{minipage}[b]{\linewidth}\raggedright
Address range
\end{minipage} & \begin{minipage}[b]{\linewidth}\raggedright
Family
\end{minipage} & \begin{minipage}[b]{\linewidth}\raggedright
Representative operations
\end{minipage} & \begin{minipage}[b]{\linewidth}\raggedright
Status
\end{minipage} \\
\midrule\noalign{}
\endhead
\bottomrule\noalign{}
\endlastfoot
0x0100--0x0106 & Inference (\texttt{inference.rs}) & deploy, inference,
batch, metadata, proof-verify, benchmark, encryption; 0x0101/0x0102
attestation-gated & Implemented \\
0x0107--0x0109 & Verification (\texttt{verify.rs}) & 0x0107 Poseidon
TENSOR\_COMMIT, 0x0108 Halo2-KZG INFERENCE\_PROOF\_VERIFY, 0x0109
MERKLE\_VERIFY\_TENSOR & Implemented (0x0108 feature-gated) \\
0x010A--0x010F & Deterministic Q16.16 compute (\texttt{compute.rs}) &
MATMUL, DOT, SOFTMAX, RELU, LINEAR, TRANSPOSE & Implemented \\
0x0110--0x0111 & Learning (\texttt{q16/belnap.rs},
\texttt{q16/routing.rs}) & Belnap 4-valued aggregation, routing &
Implemented \\
0x0130 & Recursive-fold CommD verifier (\texttt{commd\_fold\_verify.rs})
& proof-of-storage fold verification & Implemented \\
0x0200--0x0209 & x402 payments (\texttt{x402.rs}) & authorization
verify, EIP-712 verify & Implemented \\
\end{longtable}
}

Gas costs are calibrated to computational cost and include
reduced-circuit versions for the proof-verifying precompiles
(\texttt{verify.rs:89-111}).

\subsubsection{3.3 Verifiable Inference}\label{verifiable-inference}

\textbf{{[}Implemented: signature and optimistic tiers; Feature-gated:
ZK tier.{]}} Inference results carry a cost-security tradeoff across
three tiers. Signature-based verification relies on the provider's
staked identity, suitable for low-value queries where stake exceeds the
gain from cheating. Optimistic verification posts a result with a
fraud-proof challenge window; a challenger who re-executes and proves
fraud triggers slashing (Section 5; \texttt{NematocystSlashing.sol}).
ZK-SNARK verification provides cryptographic rather than economic
security: the \textbf{Halo2-KZG verifier is real code}
(\texttt{core/execution/src/zkp/halo2/mod.rs:199}, using
\texttt{verify\_proof\_multi}, \texttt{KZGCommitmentScheme},
\texttt{VerifierSHPLONK}), invoked by precompile 0x0108
(\texttt{verify.rs:118}), but it is \textbf{compiled behind the
\texttt{halo2-substrate} cargo feature}; the default build returns a
discoverable ``feature absent'' error (\texttt{mod.rs:132}). We state
this plainly: proving and verifying are implemented and tested
(\texttt{tests/inference\_proof\_verify\_real\_srs.rs},
\texttt{tests/porep\_proof\_verify\_e2e.rs}), but off by default. Paper
X, \emph{The Substrate of Verifiable Inference}, specifies this layer in
full.

\subsubsection{3.4 Model Registry and Deterministic
Compute}\label{model-registry-and-deterministic-compute}

\textbf{{[}Implemented{]}} Model registration is realized through the
inference precompile's deploy path (0x0100) together with on-chain
contracts; each registered model carries a deterministic weight hash, an
IPFS CID for decentralized weight storage, architecture metadata,
version history, and the registrant's staked identity. Deterministic
\textbf{Q16.16 fixed-point compute}
(\texttt{core/execution/src/precompiles/q16/ops.rs}) backs both the
tensor primitives (0x010A--0x010F) and the learning precompiles, with
cross-platform determinism fixtures
(\texttt{tests/cross\_platform/q16\_determinism.rs}) so that inference
used inside consensus is bit-reproducible across nodes.

\subsection{4. Orchestration Layer: Model Context
Protocol}\label{orchestration-layer-model-context-protocol}

\textbf{{[}Implemented{]}} The MCP layer (\texttt{core/mcp/}) exposes
OpenAI- and Anthropic-compatible REST endpoints, \texttt{/v1/models},
\texttt{/v1/chat/completions}, \texttt{/v1/embeddings},
\texttt{/v1/messages}, and \texttt{/v1/jobs} for async work, alongside
the JSON-RPC \texttt{eth\_*} surface at port 8545 and custom
\texttt{citrate\_*} DAG queries. The February 2026 draft listed MCP as
\emph{Specified}; the compatible REST surface is now implemented,
letting contracts and off-chain clients discover models and route
inference through a standard interface. A full decentralized routing
marketplace, in which hosts set inference prices and consumers select
providers by capability, load, and cost, remains
\textbf{{[}Specified{]}} and depends on the marketplace contracts.

\subsection{5. SALT Tokenomics}\label{salt-tokenomics}

\textbf{{[}Implemented in contracts; economic policy Specified{]}} SALT
is Citrate's native token, fixed supply 1 billion, functioning as gas,
staking collateral, and governance weight. The economic policy below is
design; the mechanisms enforcing it are deployed contracts on chain
40204.

\textbf{Table 2. SALT distribution (design).}

{\def\LTcaptype{none} % do not increment counter
\begin{longtable}[]{@{}
  >{\raggedright\arraybackslash}p{(\linewidth - 6\tabcolsep) * \real{0.2750}}
  >{\raggedleft\arraybackslash}p{(\linewidth - 6\tabcolsep) * \real{0.3000}}
  >{\raggedleft\arraybackslash}p{(\linewidth - 6\tabcolsep) * \real{0.2000}}
  >{\raggedright\arraybackslash}p{(\linewidth - 6\tabcolsep) * \real{0.2250}}@{}}
\toprule\noalign{}
\begin{minipage}[b]{\linewidth}\raggedright
Allocation
\end{minipage} & \begin{minipage}[b]{\linewidth}\raggedleft
Percentage
\end{minipage} & \begin{minipage}[b]{\linewidth}\raggedleft
Amount
\end{minipage} & \begin{minipage}[b]{\linewidth}\raggedright
Vesting
\end{minipage} \\
\midrule\noalign{}
\endhead
\bottomrule\noalign{}
\endlastfoot
Mining rewards & 50\% & 500,000,000 & emitted over network lifetime via
block rewards \\
Ecosystem development & 25\% & 250,000,000 & grants, partnerships,
developer incentives \\
Treasury & 10\% & 100,000,000 & DAO-governed reserve \\
Team & 15\% & 150,000,000 & 4-year vesting, 1-year cliff \\
\end{longtable}
}

Validators that host registered models and serve verifiable inference
earn bonuses beyond the base block reward, aligning validation with
AI-infrastructure provision. The enforcing contracts are live:
\textbf{market-maker allocation} (\texttt{MarketMakerAllocation.sol},
\texttt{0xa87fae5c…}), \textbf{contribution accounting}
(\texttt{ContributionAccounting.sol}, \texttt{0xcdd24773…}, Paper VII),
\textbf{treasury governance} (\texttt{TreasuryGovernor.sol},
\texttt{0x62e268f2…}, Paper VIII), and \textbf{tiered Byzantine
slashing} (\texttt{NematocystSlashing.sol}, \texttt{0xfeb23abd…}, Paper
IX), which implements stake, unstake, slash, an \texttt{isSlashable}
gate, and a correlated-failure multiplier. Validator membership and
staking are enforced by \texttt{CitrateMemberSBT.sol}
(\texttt{0xAD826D04…}) and \texttt{MembershipStakeVault.sol}
(\texttt{0x04c32967…}). In keeping with Citrate Inc.'s
public-communication policy, this paper makes no claim of day-one cash
earnings from compute; SALT's availability on a licensed exchange is
pending.

\subsection{6. Security Analysis}\label{security-analysis}

\subsubsection{6.1 Consensus Security}\label{consensus-security}

The GhostDAG + BFT consensus inherits well-studied properties. Safety
(no conflicting finalized states) holds under the Byzantine assumption f
\textless{} n/3 {[}11{]}; with a 100-validator committee requiring 67
signatures, this tolerates up to 33 Byzantine validators. Liveness holds
under GhostDAG's properties when fewer than half of validators are
Byzantine {[}1, 2{]}. The PoW component provides Sybil resistance for
block production while the BFT finality layer provides fast
deterministic finality, avoiding the finality delays of pure PoW while
retaining permissionless participation.

\subsubsection{6.2 Inference Verification
Security}\label{inference-verification-security}

The three-tier verification system is defense in depth. The signature
tier is bounded by the provider's stake, an explicit and auditable
economic bound. The optimistic tier depends on an honest challenger
being online, mitigated by awarding a share of the slashed stake to the
challenger. The ZK tier, when its feature is enabled, provides
cryptographic rather than economic security at higher gas cost.
\textbf{Open problem:} ZK proof generation for large-model inference
remains expensive; the practical envelope today is small models and
high-value queries. This is why 0x0108 ships feature-gated rather than
always-on, and why Paper X treats TEE attestation
(\texttt{TEEAttestationRegistry.sol}, \texttt{0x4df26aae…}) as a
complementary trust path.

\subsubsection{6.3 Dual Cryptography}\label{dual-cryptography}

\textbf{{[}Implemented{]}} Each account maintains ECDSA (secp256k1) and
Ed25519 key pairs. Ed25519 gives faster verification and deterministic
signatures, eliminating the k-nonce vulnerability that has caused ECDSA
key exposure in production systems. Hedera has operated a similar
dual-cryptography design in production since 2019, demonstrating
viability at scale.

\subsection{7. Comparative Position}\label{comparative-position}

Citrate sits at the intersection of DAG-based consensus, EVM-compatible
execution, and decentralized AI. We compare along the dimension most
relevant to each neighbor rather than forcing one table across
incommensurable systems.

\textbf{7.1 DAG consensus: Kaspa.} Kaspa is the production reference for
GhostDAG. After Crescendo (May 2025) it runs at 10 BPS with k=124 and
max parents 16, reaching a record \textasciitilde3,585 TPS {[}30{]}.
Citrate runs more conservatively (2-second blocks, k=18, max parents 10)
to carry AI-augmented block overhead, and adds a general-purpose
execution layer with AI precompiles atop the same consensus family,
which Kaspa does not provide.

\textbf{7.2 EVM chains and L2s.} Against Arbitrum, Optimism, zkSync,
Avalanche C-Chain, Moonbeam, and Hedera, Citrate's differentiator is the
AI precompile suite: no existing EVM-compatible chain provides native
precompiles for model registration, verifiable inference, deterministic
tensor compute, and paraconsistent aggregation. Applications needing
on-chain AI otherwise implement it as expensive Solidity or trust
off-chain oracles.

\textbf{7.3 Decentralized AI: Bittensor and the AI-native L1 category.}
Bittensor is an AI coordination marketplace scoring miner outputs via
Yuma consensus; Citrate is a general blockchain with composable AI
operations. Comparing them on TPS is a category error. The meaningful
comparison is integration approach: Bittensor isolates AI tasks in
subnets; Citrate embeds AI operations in the execution layer, composable
with arbitrary contract logic. Newer entrants sharpen the contrast.
Proof-of-logits designs (Ambient) make inference itself the mining
function; optimistic and TEE-attested inference rollups {[}32{]} verify
generative inference off the hot path; the DeAI systematization {[}31{]}
maps the category. Citrate's position is that verifiable inference
should be a precompile a DeFi risk engine or a model-governed DAO can
call inline, not a separate network to bridge to.

\subsection{8. Open Problems and Future
Work}\label{open-problems-and-future-work}

\textbf{8.1 Federated learning at scale.} The learning path of Paper II
is implemented at the execution and daemon layers (Belnap precompile
0x0110; \texttt{learning-daemon/src/aggregator.rs};
\texttt{LearningCheckpoint} at \texttt{learning/src/checkpoint.rs:14}).
Open questions remain on the minimum viable meta-model at
checkpoint-interval timescales, how LoRA-adapter composition scales with
accumulated adapters, and provable convergence of the recursive learning
loop under Byzantine conditions.

\textbf{8.2 Throughput scaling.} Kaspa's trajectory from 1 to 10 BPS
suggests GhostDAG systems can raise block rates substantially. Citrate's
2-second cadence is deliberately conservative for AI-augmented blocks;
with better embedding compression the rate can rise via governance,
proportionally increasing throughput. This is an engineering
optimization, not an architectural limit.

\textbf{8.3 Hardware optimization.} Paper V (ATIS) proposes analog
middleware for energy-efficient attention pruning on FPAA hardware. If
realized, it would optimize the deterministic tensor precompiles; it
remains a research proposal with no prototype built.

\textbf{8.4 Cross-chain value transfer.} Paper VI describes the Memetic
Money Portal. The honest current state is that Citrate implements a
\texttt{WrappedSALT} EIP-3009 gasless-transfer path
(\texttt{0xaa918302…}) plus market-maker allocation and a SALT/USD feed
embedded in \texttt{ComputePricingOracle.sol} (\texttt{0xdcebd5ec…});
there is no SNAP/SALT lock-mint bridge contract, and Paper VI is being
revised to describe the mechanism as it exists rather than as a bridge.

\subsection{9. Conclusion}\label{conclusion}

Citrate is a BlockDAG network that adds AI-native execution to GhostDAG
consensus. The implemented system provides parallel block production
with dual (depth + committee-BFT) finality, EVM compatibility via REVM,
families of AI precompiles for inference, verifiable proof checking,
deterministic tensor compute and paraconsistent aggregation, dual
cryptographic identity, an MCP-compatible REST surface, and SALT
tokenomics enforced by 39 deployed contracts on chain 40204.

The distinguishing contribution is not any single component but their
combination: a general-purpose programmable blockchain with native,
composable AI operations on a consensus architecture that carries
federated learning as a first-class path. Whether that combination
produces capabilities beyond its parts is an empirical question the
companion papers address theoretically (II--III), from a hardware
perspective (V), economically (VI--VII), through governance (VIII), and
through the biological design lens (IX); Paper X specifies the
verifiable-inference substrate and Paper XI presents the
Neuroarchitectural Transformer that runs on it. We have kept the earlier
draft's discipline of separating implemented from specified and
labelling targets as targets, on the conviction that credibility, not
headline numbers, is what builds lasting systems.

\subsection{Acknowledgments}\label{acknowledgments}

Drafting and literature triage were assisted by AI systems; all
mechanical claims were verified by the authors against the referenced
source files and deployed contract addresses on chain 40204. This work
received no external funding.

\subsection{References}\label{references}

{[}1{]} Sompolinsky, Y., \& Zohar, A. (2015). Secure high-rate
transaction processing in Bitcoin. \emph{Financial Cryptography and Data
Security}, 507--527. {[}2{]} Sompolinsky, Y., \& Zohar, A. (2018).
PHANTOM and GHOSTDAG: A scalable generalization of Nakamoto consensus.
\emph{IACR Cryptology ePrint Archive.} {[}3{]} Keidar, I.,
Kokoris-Kogias, E., Naor, O., \& Spiegelman, A. (2021). All you need is
DAG. \emph{Proc. ACM PODC}, 165--175. {[}4{]} Danezis, G.,
Kokoris-Kogias, L., Sonnino, A., \& Spiegelman, A. (2022). Narwhal and
Tusk. \emph{Proc. EuroSys.} {[}5{]} Spiegelman, A., Giridharan, N.,
Sonnino, A., \& Kokoris-Kogias, L. (2022). Bullshark: DAG BFT protocols
made practical. \emph{Proc. CCS.} {[}6{]} McMahan, B., et al.~(2017).
Communication-efficient learning of deep networks from decentralized
data. \emph{AISTATS}, 1273--1282. {[}7{]} Shazeer, N., et al.~(2017).
Outrageously large neural networks: the sparsely-gated
mixture-of-experts layer. \emph{ICLR.} {[}8{]} Hu, E. J., et al.~(2021).
LoRA: Low-Rank Adaptation of Large Language Models. \emph{ICLR 2022.}
{[}9{]} Vaswani, A., et al.~(2017). Attention is all you need.
\emph{NeurIPS 30.} {[}10{]} Buterin, V. (2014). Ethereum: a
next-generation smart contract and decentralized application platform.
{[}11{]} Castro, M., \& Liskov, B. (1999). Practical Byzantine fault
tolerance. \emph{OSDI}, 173--186. {[}12{]} Ben-Sasson, E., et
al.~(2014). Succinct non-interactive zero knowledge for a von Neumann
architecture. \emph{USENIX Security}, 781--796. {[}13{]} Rao, J., \&
Opentensor Foundation. (2021). Bittensor: a peer-to-peer intelligence
market. {[}14{]} Belnap, N. D. (1977). A useful four-valued logic. In
Dunn, J. M., Epstein, G. (eds), \emph{Modern Uses of Multiple-Valued
Logic}, Episteme vol.~2, Reidel, Dordrecht. {[}15{]} Kirkpatrick, J., et
al.~(2017). Overcoming catastrophic forgetting in neural networks.
\emph{PNAS}, 114(13), 3521--3526. {[}16{]} Nakamoto, S. (2008). Bitcoin:
a peer-to-peer electronic cash system. {[}17{]} Ilharco, G., et
al.~(2023). Editing models with task arithmetic. \emph{ICLR.} {[}18{]}
Yadav, P., et al.~(2023). TIES-Merging: resolving interference when
merging models. \emph{NeurIPS.} {[}19{]} Biderman, D., et al.~(2024).
LoRA learns less and forgets less. \emph{Transactions on Machine
Learning Research.} (First author Dan Biderman; corrected from the
February 2026 draft, which listed ``Biderman, S.'') {[}20{]} Weissbourd,
B., et al.~(2021). A genetically tractable jellyfish model for systems
and evolutionary neuroscience. \emph{Cell}, 184(24), 5854--5868.
{[}21{]} Pallasdies, F., et al.~(2019). From single neurons to behavior
in the jellyfish \emph{Aurelia aurita}. \emph{eLife}, 8, e50084.
{[}22{]} Klosowski, L. (2023). Mentor/Mentee Relativity: Organizational
Learning in Mentorship-Driven Swarms. Unpublished internal note, Citrate
Inc. {[}23{]} Anthropic. (2024). Model Context Protocol Specification.
Technical report. {[}27{]} Yin, M., et al.~(2019). HotStuff: BFT
consensus with linearity and responsiveness. \emph{PODC}, 347--356.
{[}28{]} Priest, G. (2006). \emph{In Contradiction: A Study of the
Transconsistent.} Oxford University Press. {[}30{]} Kaspa Network.
(2025). KIP-14: Crescendo Hardfork. Activated May 5, 2025. {[}31{]}
Zhang, W., et al.~(2024). SoK: Blockchain-Based Decentralized AI.
arXiv:2411.17461. {[}32{]} Authors TBD. (2025). Optimistic TEE-Rollups:
A Hybrid Architecture for Scalable and Verifiable Generative AI
Inference on Blockchain. arXiv:2512.20176. \emph{(Verify author list
before final submission.)}

\subsection{Appendix A: Protocol
Parameters}\label{appendix-a-protocol-parameters}

\textbf{Table A1. Citrate protocol configuration (reconciled against
code, Aug 2026).}

{\def\LTcaptype{none} % do not increment counter
\begin{longtable}[]{@{}
  >{\raggedright\arraybackslash}p{(\linewidth - 4\tabcolsep) * \real{0.4231}}
  >{\raggedright\arraybackslash}p{(\linewidth - 4\tabcolsep) * \real{0.2692}}
  >{\raggedright\arraybackslash}p{(\linewidth - 4\tabcolsep) * \real{0.3077}}@{}}
\toprule\noalign{}
\begin{minipage}[b]{\linewidth}\raggedright
Parameter
\end{minipage} & \begin{minipage}[b]{\linewidth}\raggedright
Value
\end{minipage} & \begin{minipage}[b]{\linewidth}\raggedright
Source
\end{minipage} \\
\midrule\noalign{}
\endhead
\bottomrule\noalign{}
\endlastfoot
Chain ID & 40204 (testnet) & \texttt{testnet-config.toml:2};
\texttt{genesis\_model.rs:269} \\
Block time & \textasciitilde2 s (testnet), 1 s (devnet) &
\texttt{testnet-config.toml:4}; \texttt{block\_builder.rs:76} \\
Max parents & 10 (1 selected + 9 merge) &
\texttt{core/consensus/src/types.rs:179} \\
k parameter & 18 & \texttt{core/consensus/src/types.rs:178} \\
Checkpoint interval & 50 blocks &
\texttt{core/consensus/src/checkpoint.rs:101} \\
Committee size & 100 validators & \texttt{checkpoint.rs} \\
BFT quorum & 67 (2/3+1) & \texttt{checkpoint.rs:88,103} \\
Vote domain separator & \texttt{CITRATE-CHECKPOINT-V1} &
\texttt{checkpoint.rs:59} \\
Signature schemes & ECDSA (secp256k1) + Ed25519 &
\texttt{core/execution/src/types.rs} \\
Deployed contracts (40204) & 39 &
\texttt{contracts/DEPLOYED\_ADDRESSES.md} \\
Core test count & 3,517 & \texttt{cargo\ test\ -\/-workspace} (grep) \\
Parallel-execution gate & 8-worker ≥ 2× 1-worker &
\texttt{execution/benches/tps\_parallel.rs:10} \\
TPS figure (target, not benched) & 5,000 sustained / 10,000 ceiling &
\texttt{README.md:65} \\
SALT total supply & 1,000,000,000 & design \\
\end{longtable}
}

\begin{center}\rule{0.5\linewidth}{0.5pt}\end{center}

\emph{This paper is part of the Gradient Papers series, published by
Citrate Inc.} \emph{Correspondence: Larry@citrate.ai}
